%%%%%%%%%%%%%%%%%%%%%%%%%%%%%%%%%%%%%%%%%%%%%%%%%%%%%%%%%%%%%%%%%%%%%%%%%%%%%%%
%% Quick Reference Guide (QRF) Infrastructure                                %%
%%%%%%%%%%%%%%%%%%%%%%%%%%%%%%%%%%%%%%%%%%%%%%%%%%%%%%%%%%%%%%%%%%%%%%%%%%%%%%%

% --- Compact QRF Main Title ---
\NewDocumentCommand{\qrfTitle}{m}{%
  \clearpage           % Start on a fresh page
  \phantomsection      % Anchor for hyperref/TOC
  \addcontentsline{toc}{chapter}{#1} % Add to Table of Contents
  \markboth{#1}{#1}    % Set the header on the left/right pages
  
  \begin{center}
      \Large \bfseries #1 % Compact, centered title
  \end{center}
  \vspace{-1ex} % Pull the columns up slightly closer to the title
}

% --- 1. Configurable Parameters ---

% Column and Equation Spacing (Absolute lengths)
\newlength{\qrfColumnSep}
\setlength{\qrfColumnSep}{0.5cm}       % Gap between the two columns
\newlength{\qrfEqSkip}
\setlength{\qrfEqSkip}{4pt}          % Vertical space above/below equations
\newlength{\qrfHeaderPadding}
\setlength{\qrfHeaderPadding}{1em}   % Space between the header text and the whiskers

% Header Whisker & Vertical Spacing (Font-relative sizes)
\newcommand{\qrfHeaderSpaceAbove}{1.5ex}
\newcommand{\qrfHeaderSpaceBelow}{1ex}
\newcommand{\qrfWhiskerHeight}{0.65ex}
\newcommand{\qrfWhiskerDepth}{-0.55ex}


% --- 2. Counters ---
\newcounter{qrfEq}
\renewcommand{\theqrfEq}{Q.\arabic{qrfEq}}
\renewcommand{\theHqrfEq}{Q.\arabic{qrfEq}} % Prevents hyperref collisions


% --- 3. Macros ---

% Usage: \qrfeq[label]{math}{citation}
\NewDocumentCommand{\qrfeq}{o m m}{%
  \refstepcounter{qrfEq}%
  \IfValueT{#1}{\label{#1}}%
  \begin{equation*}
    #2 \tag*{(\theqrfEq)\quad{\footnotesize #3}}
  \end{equation*}%
}

% Centered QRF header with dynamic "whisker" lines
\NewDocumentCommand{\qrfheader}{m}{%
  \par\vspace{\qrfHeaderSpaceAbove}
  \noindent
  \leavevmode
  % Left whisker
  \leaders\hrule height \qrfWhiskerHeight depth \qrfWhiskerDepth\hfill
  % Centered Text
  \hspace{\qrfHeaderPadding} \textbf{\large #1} \hspace{\qrfHeaderPadding}
  % Right whisker
  \leaders\hrule height \qrfWhiskerHeight depth \qrfWhiskerDepth\hfill
  \null
  \par\vspace{\qrfHeaderSpaceBelow}
}


% --- 4. The Main Layout Environment ---
% This safely applies the column widths and equation spacing ONLY to the QRF
\newenvironment{qrfLayout}{%
    \setlength{\columnsep}{\qrfColumnSep}
    % \setlength{\columnseprule}{0.4pt} % Uncomment for a vertical divider line!
    \begin{multicols}{2}
    \setlength{\abovedisplayskip}{\qrfEqSkip}
    \setlength{\belowdisplayskip}{\qrfEqSkip}
}{%
    \end{multicols}
}




\qrfTitle{Quick Reference Guide}
\addcontentsline{toc}{chapter}{Quick Reference Guide}
\markboth{Quick Reference Guide}{Quick Reference Guide}

\begin{qrfLayout}

\qrfheader{Commutation Relations}

\qrfeq[qrf:momentum]{
    [x, p] = i\hbar
}{\cite{1.19:wikipediaVirial}}

\qrfeq{
    [S_x, S_y] = i\hbar S_z
}{[Prob \inref{some_label}]}
\qrfeq[qrf:momentum]{
    [x, p] = i\hbar
}{\cite{1.19:wikipediaVirial}}

\qrfeq{
    [S_x, S_y] = i\hbar S_z
}{[Prob \inref{some_label}]}
\qrfeq[qrf:momentum]{
    [x, p] = i\hbar
}{\cite{1.19:wikipediaVirial}}

\qrfeq{
    [S_x, S_y] = i\hbar S_z
}{[Prob \inref{some_label}]}


\qrfheader{Uncertainty Principles}

\qrfeq[qrf:generalized_uncertainty]{
    \Delta A \Delta B \ge \frac{1}{2} \abs{\ev{[A, B]}}
}{\cite{1.20:wikipedia:StrongerUncertaintyRelations}}
\qrfeq[qrf:generalized_uncertainty]{
    \Delta A \Delta B \ge \frac{1}{2} \abs{\ev{[A, B]}}
}{\cite{1.20:wikipedia:StrongerUncertaintyRelations}}
\qrfeq[qrf:generalized_uncertainty]{
    \Delta A \Delta B \ge \frac{1}{2} \abs{\ev{[A, B]}}
}{\cite{1.20:wikipedia:StrongerUncertaintyRelations}}
\qrfeq[qrf:generalized_uncertainty]{
    \Delta A \Delta B \ge \frac{1}{2} \abs{\ev{[A, B]}}
}{\cite{1.20:wikipedia:StrongerUncertaintyRelations}}
\qrfeq[qrf:generalized_uncertainty]{
    \Delta A \Delta B \ge \frac{1}{2} \abs{\ev{[A, B]}}
}{\cite{1.20:wikipedia:StrongerUncertaintyRelations}}
\qrfeq[qrf:generalized_uncertainty]{
    \Delta A \Delta B \ge \frac{1}{2} \abs{\ev{[A, B]}}
}{\cite{1.20:wikipedia:StrongerUncertaintyRelations}}
\qrfeq[qrf:foo]{
    \Delta A \Delta B \ge \frac{1}{2} \abs{\ev{[A, B]}}
}{\cite{1.20:wikipedia:StrongerUncertaintyRelations}}

\end{qrfLayout}




%% % Configurable spacing for the Quick Reference columns
%% \newlength{\qrfColumnSep}
%% \setlength{\qrfColumnSep}{1cm}

%% % --- Quick Reference Counters & Macros ---
%% \newcounter{qrfEq}
%% \renewcommand{\theqrfEq}{Q.\arabic{qrfEq}}
%% \renewcommand{\theHqrfEq}{Q.\arabic{qrfEq}} % Prevents hyperref anchor collisions

%% % Usage: \qrfeq[label]{math}{citation}
%% \NewDocumentCommand{\qrfeq}{o m m}{%
%%   % Step the counter and drop the label outside the equation
%%   % so that \ref{} strictly captures "Q.n" and not the citation text.
%%   \refstepcounter{qrfEq}%
%%   \IfValueT{#1}{\label{#1}}%
%%   \begin{equation*}
%%     #2 \tag*{(\theqrfEq)\quad{\footnotesize #3}}
%%   \end{equation*}%
%% }

%% % Centered QRF header with "whisker" lines
%% \NewDocumentCommand{\qrfheader}{m}{%
%%   \par\vspace{1.5ex} % Space above
%%   \noindent
%%   \leavevmode
%%   % Left whisker
%%   \leaders\hrule height 0.65ex depth -0.55ex\hfill
%%   % Centered Text
%%   \quad \textbf{\large #1} \quad 
%%   % Right whisker
%%   \leaders\hrule height 0.65ex depth -0.55ex\hfill
%%   %
%%   \null % <--- THIS IS THE MAGIC FIX
%%   \par\vspace{1ex} % Space below
%% }

%% \chapter*{Quick Reference Guide}
%% \addcontentsline{toc}{chapter}{Quick Reference Guide}
%% \markboth{Quick Reference Guide}{Quick Reference Guide}

%% % Reduces vertical whitespace to pack the two columns tightly
%% \setlength{\abovedisplayskip}{4pt}
%% \setlength{\belowdisplayskip}{4pt}

%% \begin{multicols}{2}

%% \qrfheader{Commutation Relations}

%% % Example with a label! You can now use \ref{qrf:momentum} anywhere else
%% % in your manual and it will correctly print "Q.1"
%% \qrfeq[qrf:momentum]{
%%     [x, p] = i\hbar
%% }{\cite{1.19:wikipediaVirial}}

%% % Example without a label
%% \qrfeq{
%%     [S_x, S_y] = i\hbar S_z
%% }{[Prob \inref{some_label}]}

%% \qrfeq[qrf:product_rule]{
%%   [A,BC] = [A,B]C + B[A,C]
%% }{[Prob \inref{another_label}]}


%% \qrfheader{Uncertainty Principles}

%% \qrfeq[qrf:generalized_uncertainty]{
%%     \Delta A \Delta B \ge \frac{1}{2} \abs{\ev{[A, B]}}
%% }{\cite{1.20:wikipedia:StrongerUncertaintyRelations}}

%% \end{multicols}

%% % Restore normal spacing for the rest of the book
%% \setlength{\abovedisplayskip}{10pt}
%% \setlength{\belowdisplayskip}{10pt}


%% \renewcommand{\thechapter}{Q}
%% \chapter{Quick Reference}

%% Because the $'$ and $''$ and $'''$ and so on gets really annoying, I
%% switch from primed Eigenvectors/Eigenvalues to subscripts.  So you'll
%% see $\ket{a'} \Rightarrow \ket{a_i}$ and $\ket{a''} \Rightarrow
%% \ket{a_j}$ and so on.

%% \begin{equation}
%%   \qlabel{taylor:exponential}
%%   e^{A} \equiv \sum_{k=0}^{\infty} \frac{1}{k!} A^k
%% \end{equation}

%% Not sure about including this one.
%% \begin{equation}
%%   \inlabel{taylor:rotation}
%%   \begin{split}
%%     \exp(\frac{i\sigma_z\phi}{2})
%%     &= \sum_{0}^{\infty}{\frac{1}{k!}\left(\frac{i\sigma_z\phi}{2}\right)^k}
%%     \\
%%     &= \cos(\frac \phi 2) + i \sigma_z \sin(\frac \phi 2)
%%     \\
%%   \end{split}
%% \end{equation}\outref{1}{5}{fixme}

%% \section{Pauli Spin}

%% \begin{equation*}
%%   \tag{1.110a}
%%   \qlabel{spin:x}
%%   \ket{x\pm} = \frac{1}{\sqrt{2}} [\ket + \pm \ket -]
%% \end{equation*}

%% \begin{equation*}
%%   \tag{1.110b}
%%   \qlabel{spin:y}
%%   \ket{y\pm} = \frac{1}{\sqrt{2}} [\ket + \pm i\ket -]
%% \end{equation*}

%% \begin{equation*}
%%   \tag{1.111a}
%%   \qlabel{spin:sx}
%%   S_x = \frac{\hbar}{2} [\ket - \bra + + \ket + \bra - ]
%% \end{equation*}

%% \begin{equation*}
%%   \tag{1.111b}
%%   \qlabel{spin:sy}
%%   S_y = i \frac{\hbar}{2} [\ket - \bra + - \ket + \bra - ]
%% \end{equation*}

%% \begin{equation*}
%%   \tag{3.50}
%%   \qlabel{spin:matrices}
%%   \sigma_x =
%%   \begin{pmatrix}
%%     0 & 1 \\
%%     1 & 0 \\
%%   \end{pmatrix},
%%   \sigma_y =
%%   \begin{pmatrix}
%%     0 & -i \\
%%     i &  0 \\
%%   \end{pmatrix},
%%   \sigma_z =
%%   \begin{pmatrix}
%%     1 &  0 \\
%%     0 & -1 \\
%%   \end{pmatrix}
%% \end{equation*}

%% \begin{figure}
%%   \begin{center}
%%     \qlabel{spinAxes}
%%     \includegraphics[width=2in]{qrf/spin_axes.png}
%%     \caption{Spin axis definitions.  Please note that due to font
%%       limitations in QCad at the time of this writing, $\alpha$ and
%%       $\beta$ are represented as $a$ and $b$ respectively.}
%%   \end{center}
%% \end{figure}

%% \begin{equation}
%%   \qlabel{spinAngle}
%%   \ket{\alpha} =
%%   \cos(\frac{\beta}{2})\ket{+} + \sin(\frac{\beta}{2}) e^{i\alpha}\ket{-}
%% \end{equation}\outref{1}{11}{spinAngle}

%% \subsection{Integral Table}

%% \begin{equation}
%%   \qlabel{integral:gaussian}
%%   \int_{-\infty}^{\infty} \dd{x} e^{-\alpha x^2} = \sqrt{\frac{\pi}{\alpha}}
%% \end{equation}\cite[eq 3.321.3]{gr}

%% \begin{equation}
%%   \qlabel{uncertain:full}
%% \dispersion{A}\dispersion{B} \ge \frac{1}{4} \abs{\expval{[A, B]}}^2 + \frac{1}{4} \abs{\expval{{A, B} - 2\expval{A}\expval{B}}}^2
%% \end{equation}\outref{1}{20}{theQuestion}

%% \begin{equation}
%%   \qlabel{integral:xgaussian}
%%   \int_{-\infty}^{\infty} \dd{x} xe^{-\alpha x^2} = 0
%% \end{equation}\cite[eq 3.321.3]{gr}

%% \begin{equation}
%%   \qlabel{integral:xxgaussian}
%%   \int_{-\infty}^{\infty} \dd{x} x^2e^{-\alpha x^2} = \frac{1}{2\alpha}\sqrt{\frac{\pi}{\alpha}}
%% \end{equation}\cite[fixme]{gr}

%% \section{Wave Functions}

%% \subsection{\acf{gwp}}
%% \qlabel{wavey:gwp}
%%     \begin{equation}
%%       \qlabel{gwp:exp}
%%       \begin{alignedat}{2}
%%         \expval{X}^2 &= 0 & \quad{} &\text{(eq \ref{1.24:exp:x1})}\\
%%         \expval{X^2} &= \frac{d^2}{2} & &\text{(eq \ref{1.24:exp:x2})}\\
%%         \expval{P}^2 &= \hbar^2k^2 & &\text{(eq \ref{1.24:exp:p1})}\\
%%         \expval{P^2} &= \hbar^2\left(k^2+\frac{1}{2d^2}\right)
%%         & &\text{(eq \ref{1.24:exp:p2})}\\
%%       \end{alignedat}
%%     \end{equation}

%%     \begin{equation}
%%       \qlabel{uncertain:xp}
%%       \dispersion{X}\dispersion{P}
%%       =
%%       \left(\expval{\hat{X}^2} - \expval{\hat{X}}^2\right)
%%       \left(\expval{\hat{P}^2} - \expval{\hat{P}}^2\right)
%%       = \frac{\hbar^2}{4}
%%     \end{equation}

%%     \begin{equation}
%%       \qlabel{uncertain:et}
%%       \dispersion{E}\dispersion{T}
%%       =
%%       FIXME
%%     \end{equation}FIXME: Cite wikipedia like I did in problem 24

%%     \begin{equation}
%%       \qlabel{wavey:gwp:psi}
%%       \psi(x')
%%       = \braket{x'}{\alpha}
%%       = \frac{1}{\pi^{1/4}}\sqrt{\frac{1}{d}}
%%       \exp[ikx'-\frac{x'^2}{2d^2}]
%%     \end{equation}

%%     \begin{equation}
%%       \qlabel{wavey:gwp:phi}
%%       \phi(p')
%%       =
%%       \left[ \frac{1}{\sqrt{2\pi\hbar}} \right]
%%       \int \dd{x'} \exp(\frac{-ip'x'}{\hbar})\psi_{\alpha}(x')
%%       = 
%%       \frac{1}{\pi^{1/4}}
%%       \sqrt{\frac{d}{\hbar}}
%%       \exp[-\frac{d^2}{2\hbar^2}(p' - \hbar k)^2]
%%     \end{equation}

%%     \begin{equation}
%%       \qlabel{comms:chain}
%%       \tag{FIXME}
%%       [A,BC] = [A,B]C + B[A,C]
%%     \end{equation}

%%     \begin{equation}
%%       \qlabel{comms:power}
%%       \begin{split}
%%         [A,B^n] &= n B^{n-1}[A,B]\\
%%         [A,B] &= C(\text{onstant})\\
%%       \end{split}
%%     \end{equation}\outref{1}{30}{ConstantPower}
    
%%     \begin{equation}
%%       \qlabel{comms:func}
%%       [x, F(p)] = i\hbar \pdv{F}{p}, \quad [p, F(x)] = -i\hbar \pdv{F}{x}
%%     \end{equation}\outref{1}{30}{xfp}

%%     \begin{equation}
%%       \qlabel{misc:realeig}
%%       A = A^{\dagger} \implies \forall i, \lambda_i \in \mathscr{R}
%%     \end{equation}

%%     \begin{equation}
%%       \qlabel{misc:unitary}
%%       U^{\dagger}U = I \implies \text{U is Unitary}
%%     \end{equation}

%%     \begin{equation}
%%       \qlabel{decomp:func}
%%       f(A) = \int \dd{a'} f(a') \Lambda_{a'}
%%     \end{equation}\outref{1}{29}{mainPoint}

%%     \begin{equation}
%%       \qlabel{decomp:ident}
%%       I = \sum_i \Lambda_i \implies I = \int \dd{a'} \Lambda_{a'}
%%     \end{equation}

%%   \begin{equation}
%%       \qlabel{op:p}
%%     \mel{x'}{\hat{p}}{\alpha} = -i\hbar\pdv{}{x'}\braket{x'}{\alpha}
%%   \end{equation}

%%   \begin{equation}
%%       \qlabel{op:x}
%%     \mel{p'}{\hat{x}}{\alpha} = i\hbar\pdv{}{p'}\braket{p'}{\alpha}
%%   \end{equation}

%%   \begin{equation}
%%       \qlabel{stats:var:minus}
%%     \dispersion{A} = \expval{A^2} - \expval{A}^2
%%   \end{equation}

%%   \begin{equation}
%%       \qlabel{stats:var:func}
%%     \dispersion{f(X)} \approx \left(f'(\expval{X})\right)^2 \dispersion{X}
%%   \end{equation}


%% \subsection{Trig Identities}

%% \begin{equation}
%%       \qlabel{trig:sin2a}
%%   \sin{2\alpha} = 2\sin{\alpha}\cos{\alpha}
%% \end{equation}

%% \begin{equation}
%%   \qlabel{trig:cos2a}
%%   \begin{split}
%%     \qlabel{trigCos2A}
%%     \cos{2\alpha}
%%     &= \cos^2{\alpha} - \sin^2{\alpha}\\
%%     &= 2\cos^2{\alpha} - 1\\
%%     &= 1 - 2\sin^2{\alpha}\\
%%   \end{split}
%% \end{equation}

%% \begin{equation}
%%   \qlabel{trig:coshalf}
%%   \cos{\frac \alpha 2} = \pm \sqrt{\frac{1 + \cos{\alpha}}{2}}
%% \end{equation}

%% \begin{equation}
%%   \qlabel{trig:sinhalf}
%%   \sin{\frac \alpha 2} = \pm \sqrt{\frac{1 - \cos{\alpha}}{2}}
%% \end{equation}

%% \begin{equation}
%%   \qlabel{trig:harmonic}
%%   \begin{alignedat}{2}
%%     && f(\theta) &= a\cos(\theta) + b\sin(\theta)\\
%%     &&&= c\cos(\theta + \delta)\\
%%     \text{where } && c &= \sqrt{a^2 + b^2}\\
%%     \text{and } && \delta &= \tan^{-1}(-\frac{a}{b})
%%   \end{alignedat}
%% \end{equation}
%% \cite{msu:harmonicAddition}
